\newpage
\addcontentsline{toc}{chapter}{\MakeUppercase{Abstract}}
\thispagestyle{empty}
%
\begin{spacing}{1.2}
%
\chapter*{Abstract}

GramaGIS is a web-based Geographic Information System (GIS) platform developed for Vazhathope Grama Panchayath to digitize and centralize panchayat-level infrastructure data. It consolidates spatial and non-spatial information such as wards, roads, schools, hospitals, public offices, and citizen feedback into a single interactive web portal. The system enables administrators and citizens to visualize, manage, and access real-time geographic data through a user-friendly interface. This digital framework supports data-driven decision-making, promotes transparency, and enhances the efficiency of local governance.

Currently, most panchayat infrastructure records are maintained in paper form or stored across disconnected files, making retrieval and updates difficult. This fragmentation leads to inefficiencies in maintenance, planning, and service delivery. GramaGIS addresses these issues by integrating essential datasets into a PostgreSQL/PostGIS-backed environment and exposing them through a Leaflet.js frontend, a Node.js and Express.js application backend, GeoServer services, and a Gemini-assisted natural language query workflow. By enabling centralized geospatial management, controlled administrative updates, and public feedback collection, GramaGIS strengthens planning capabilities, simplifies public data access, and lays the foundation for a smart, spatially aware governance model. Ultimately, the platform aims to empower local bodies to make informed, transparent, and sustainable decisions for community development.

%%
\end{spacing}
